\chapter{39}



Louis kam mit Block und Schreibzeug die Treppe in die
Küche hinunter. Er schaute sich um. Wo hatte Joh nur den Ordner
hingelegt? Er wollte sich heute konzentriert in die Themen rund um
die Schafe, Alpen und Wölfe einlesen.
Er trat vors Haus und ging die wenigen Schritte hinunter in Manuelas
Richtung. Sie stand unter einem grossen Wäscheschirm im Garten, nahm
Frottiertücher von der Leine, faltete sie lose zusammen und legte sie in
eine Wäschezaine.

«Weisst du vielleicht, wohin Joh den Info-Ordner gelegt hat?»

Louis blieb vor ihr stehen. Manuela runzelte die Stirn.

«Der müsste irgendwo im Wohnzimmer liegen.»

Interessiert blickte sie auf die kleinen Kopfhörer in Louis’ Hand.

«Schön, wieder mal einer, der mit Musik arbeitet?» sie schmunzelte, «wie
ich.»

Louis gefiel ihre Bemerkung.

«Dann kannst du mir bestimmt etwas empfehlen, vielleicht vorwiegend
musikalisch ohne allzu dominanten Gesang? Der lenkt mich sonst vom Lesen
ab.»

«Also, etwas Ruhiges, hm, warte.»

Sie presste die Augen zusammen. Dann schnippte sie mit den Fingern.

«Vielleicht Michel Pépé \emph{«Le mantra du coeur»}?», sie zögerte,
«ich weiss nicht, ob dir so etwas gefällt, es ist meditativ, wirkt
beruhigend, jedenfalls auf mich. Und, ich finde es passt gut in
die Umgebung hier auf die Alp.»

«Aha, warum nicht. Mit Meditation kenn ich mich zwar nicht aus, aber
Horizonterweiterung kann nie schaden,» er machte eine unbeholfene
Grimasse, «danke, Manuela, ich werde gleich reinhören.»

Louis ging auf Ordnersuche. Er fand ihn am Boden neben dem Sofa. Dann setzte er sich
vor dem Haus in die Wiese, legte das Dossier vor sich ins Gras und
suchte auf Spotify nach Michel Pépé.
Das Stück begann mit leisen Tönen. Eine zurückhaltende Männerstimme
legte sich mitten auf einen breit ausgelegten Klangteppich. Louis schrak zusammen. Irgend
etwas stimmte da nicht. Die langgezogenen Harmonien wirkten geradezu
bedrohlich. Und der Gesang zog ihn mit einer Entschlossenheit an sich,
die ihn fürchten liess, nicht wieder freizukommen. Mit einem Mal
erschien ihm der Sänger als gefesselter, um Gnade flehender Mann. Louis
wurde übel. Ihn fror. Seine Augen begannen zu zucken. Und sein Hals
schnürte sich zu. Die Musik schien als schwarzes, flatterndes Tuch über
ihm zu schweben. Es kam näher. Würde ihn zudecken und ersticken.
Eine befremdende Traurigkeit löste den Schrecken ab. Louis starrte
entsetzt auf seine Hände, die wieder unkontrolliert zu zittern anfingen.

\qrblock{Michel Pépé \emph{«Le Mantra du Coeur»}}{media/QR-Codes/033_Le_Mantra_du_Coeur_Michel_Pépé_Spotify.png}

Louis stand auf, ging taumelnd rückwärts und beendete das Stück, bevor
er sich auf den Vorplatz des Hauses rettete.
Manuela schaute irritiert zu ihm hoch, liess ein Frottiertuch auf die
Zaine fallen und ging mit entschiedenen Schritten auf ihn zu. Mit geweiteten
Augen liess sich Louis auf die Bank fallen. Erst blieb Manuela stumm neben ihm
stehen. Dann setzte sie sich zu ihm.

«Oh,» begann sie flüsternd und legte behutsam ihre Hand auf Louis’ Arm,
«das Stück scheint dich nicht gerade beruhigt zu haben - »

«Diese Klänge - es, es ist das gleiche Gefühl,» er schluckte und atmete tief ein, «wie
damals, als ich mit, hm, ach,» mit einer forschen Handbewegung wischte er
die Erinnerung weg.

Nicht noch einmal. Kein einziges Mal mehr. Auf dieser Bank sitzen und
diese Verlorenheit fühlen.
Seinen Blick in die Ferne gerichtet, legte Louis nun seine Hände
ineinander, als halte er sich an sich selbst fest. Er
schaute Manuela mit einem müden Lächeln an.

«Wie peinlich, was ich hier abziehe, schon gestern Abend. Es tut mir
leid, Manuela, ich fühle mich gerade wie ein beknackter Idiot. Dein Tipp
war bestimmt gut gemeint.»

Was war nur in ihn gefahren, schon wieder. Warum überfiel ihn diese uralte Kindergeschichte? Warum
hier? Überhaupt, was machte dieser Ort mit ihm? 
Manuela zog ihre Lippen zu einem schmalen Lächeln hoch. Louis erkannte
darin Verständnis. Sie füllte Apfelsaft in ihre beiden Gläser, überschlug die Beine und sah aus, als stelle sie sich auf längeres Zuhören ein.
Louis blickte an sich hinunter.

«Ich will doch schwer hoffen, Ciro,» sie suchte seinen Blick, «dass wir
dir gestern keinen Anlass gaben, dich wie ein Idiot zu fühlen. Nicht nur
dein Gesang war ein Geschenk, auch wie offen du dich uns gezeigt hast.
Klar, wir kennen uns noch kaum, trotzdem, du hast Platz bei uns. Wenn du
reden willst, höre ich dir gerne zu.»

Damit legte sich Manuela die Hände über ihren Bauch und wartete. 

«Gerne, vielleicht später mal,» er versuchte, seinen Atem zu
kontrollieren, «erstmal will ich mich mit den Infos zu meiner Arbeit
beschäftigen.»

Er wusste nicht annähernd, was er ihr erzählen sollte. Dieser Gesang. Er
wollte nichts schneller, als davon loskommen.

«Ich lese mir dann mal die \emph{Anleitung zum Schäfer} durch.»

Es gelang ihm nur knapp, engagiert zu klingen.

«Ja, klar.»

Mit einem Seitenblick zu ihm stand Manuela langsam auf.

«Wenn du Fragen hast, findest du mich irgendwo im oder ums Haus, okay?»

Louis war erleichtert, dass sie nicht weiter nachfragte.

«Ja, danke, mach ich.»

Er stand ruckartig auf und setzte sich zurück auf die Wiese.